% ID: SEL-D0036-OB-E02
% ESTADO: validada. S037, cotejada con examen y criterios.
\textbf{a)} La continuidad en $-1$ exige
$-a/2-6=-1-2b$. Las derivadas laterales son $a/2$ y
$2+b$, de modo que la derivabilidad exige $a/2=2+b$.
Sustituyendo en la primera ecuación, $-8-b=-1-2b$.
Resulta $\boxed{b=7,\ a=18}$. Estos valores hacen coincidir
tanto los límites con el valor como las pendientes.

\textbf{b)} Como $-2<-1$, se usa el tramo lineal.
Con $a=1$, $f(-2)=-7$ y $f'(-2)=1/2$.
La tangente $y+7=(x+2)/2$ queda $\boxed{y=x/2-6}$.
